% ===========================================================
%  Programacao Concorrente -- Trabalho 1 -- Versao 2.0
%  Relatorio: Produtor-Consumidor com Buffer Limitado
%             e Logger Assincrono com Semaforos POSIX
% ===========================================================
\documentclass[12pt, a4paper]{article}

% ---- Pacotes disponiveis no texlive-latex-base ----
\usepackage[utf8]{inputenc}
\usepackage[T1]{fontenc}
\usepackage{geometry}
\usepackage{graphicx}
\usepackage{amsmath}
\usepackage{amssymb}
\usepackage{hyperref}
\usepackage{indentfirst}
\usepackage{array}
\usepackage{tabularx}
\usepackage{url}

% ---- Margens ----
\geometry{
    top    = 3cm,
    bottom = 2cm,
    left   = 3cm,
    right  = 2cm
}

% ---- Espacamento 1.5 manual ----
\renewcommand{\baselinestretch}{1.5}

% ---- Hyperlinks ----
\hypersetup{
    colorlinks = true,
    linkcolor  = black,
    citecolor  = black,
    urlcolor   = blue
}

% ---- Traducoes manuais (sem babel) ----
\renewcommand{\contentsname}{Sumário}
\renewcommand{\refname}{Referências}
\renewcommand{\appendixname}{Apêndice}
\renewcommand{\figurename}{Figura}
\renewcommand{\tablename}{Tabela}

% ===========================
%  Inicio do documento
% ===========================
\begin{document}
\pagenumbering{gobble}

% ---- CAPA ----
\begin{titlepage}
    \centering
    \vspace*{1cm}

    {\large \textbf{Universidade} \par}
    \vspace{0.3cm}
    {\large Departamento de Ci\^encia da Computa\c{c}\~ao \par}
    \vspace{0.3cm}
    {\large Programa\c{c}\~ao Concorrente \par}

    \vspace{3cm}

    {\Huge \textbf{Trabalho 1 --- Vers\~ao 2.0} \par}
    \vspace{0.5cm}
    {\LARGE Logger Ass\'\i{}ncrono com Sem\'aforos POSIX \par}
    \vspace{0.5cm}
    {\Large Produtor-Consumidor com Buffer Limitado e Logger Ass\'\i{}ncrono \par}

    \vspace{3cm}

    {\large \textbf{Aluno:} Gustavo Choueiri \par}
    \vspace{0.5cm}
    {\large \textbf{Data:} Novembro de 2026 \par}

    \vfill

    {\large Relat\'orio apresentado como requisito parcial da disciplina
    de Programa\c{c}\~ao Concorrente. \par}
\end{titlepage}

% ---- Sumario ----
\tableofcontents
\newpage

% ===========================
%  1. INTRODUCAO
% ===========================
\pagenumbering{arabic}
\section{Introdução}

A programação concorrente é um paradigma fundamental na computação moderna,
presente desde sistemas embarcados de baixo nível até servidores de alta escala
na nuvem. A capacidade de executar múltiplas tarefas simultaneamente aumenta
significativamente o desempenho de aplicações; contudo, introduz uma classe de
desafios delicada: a \textbf{sincronização entre processos e threads} que
compartilham recursos na memória.

Quando múltiplas threads acessam e modificam uma estrutura de dados compartilhada
sem nenhum mecanismo de controle, surgem as chamadas \textbf{condições de corrida}
(\textit{race conditions}). Uma condição de corrida ocorre quando o resultado
correto de um programa depende da ordem de execução relativa entre as threads,
tornando o comportamento da aplicação não determinístico, difícil de depurar e
potencialmente catastrófico em sistemas críticos \cite{tanenbaum}.

Para tratar esse problema, a biblioteca padrão POSIX define a interface
\texttt{Pthreads} (\textit{POSIX Threads}), que fornece primitivas robustas de
sincronização: \textbf{mutexes} (\textit{mutual exclusion locks}),
\textbf{variáveis de condição} e \textbf{semáforos POSIX} (\texttt{sem\_t}).
Essas primitivas permitem que o programador delimite regiões críticas do código,
garantindo exclusão mútua e comunicação segura entre threads \cite{kerrisk}.

O presente trabalho, em sua \textbf{Versão 2.0}, expande a solução original do
Problema Produtor-Consumidor com duas contribuições principais: (1)~a
\textbf{refatoração do código em múltiplos módulos} (\texttt{.h}/\texttt{.c}),
com responsabilidades claramente separadas; e (2)~a adição de uma
\textbf{camada de logger assíncrono} sincronizada exclusivamente com
\textbf{semáforos POSIX} (\texttt{sem\_t}), demonstrando essa primitiva em
paralelo com o mutex e as variáveis de condição já utilizados no buffer principal.
O relatório está organizado da seguinte forma: a
Seção~\ref{sec:formalizacao} formaliza o problema proposto; a
Seção~\ref{sec:algoritmo} descreve o algoritmo e a nova arquitetura; e
a Seção~\ref{sec:conclusao} apresenta as considerações finais.

% ===========================
%  2. FORMALIZACAO DO PROBLEMA
% ===========================
\section{Formalização do Problema Proposto}
\label{sec:formalizacao}

\subsection{Definição}

O \textbf{Problema Produtor-Consumidor com Buffer Limitado}
(também denominado \textit{Bounded-Buffer Problem}) é um dos problemas de
sincronização mais estudados na literatura \cite{dijkstra}. Ele modela
situações reais em que componentes de um sistema produzem dados em taxas
distintas de seu consumo, como \textit{pipelines} de processamento, filas de
mensagens e \textit{spoolers} de impressão.

Formalmente, o problema é definido por um conjunto de entidades e restrições:

\begin{itemize}
    \item \textbf{Buffer compartilhado:} uma estrutura circular de capacidade
          $N$ posições, acessível por todas as threads.
    \item \textbf{Threads Produtoras} ($P_1, P_2, \ldots, P_m$): cada uma
          gera itens e os deposita no buffer.
    \item \textbf{Threads Consumidoras} ($C_1, C_2, \ldots, C_k$): cada uma
          retira itens do buffer e os processa.
    \item \textbf{Restrição de capacidade:} o número de itens no buffer a
          qualquer instante $t$ satisfaz $0 \leq |B(t)| \leq N$.
\end{itemize}

\subsection{Condições de Corrida Identificadas}

Sem mecanismos de sincronização, as seguintes condições de corrida emergem:

\begin{enumerate}
    \item \textbf{Inserção concorrente:} dois produtores calculam simultaneamente
          o índice de cauda (\texttt{cauda}) e sobrescrevem o mesmo slot do buffer,
          perdendo um dos itens produzidos.

    \item \textbf{Remoção concorrente:} dois consumidores leem simultaneamente o
          mesmo item da cabeça (\texttt{cabeca}), processando o mesmo valor duas
          vezes e pulando o próximo.

    \item \textbf{Overflow:} um produtor insere um item quando o buffer está
          cheio ($|B(t)| = N$), corrompendo um item ainda não consumido.

    \item \textbf{Underflow:} um consumidor retira um item quando o buffer está
          vazio ($|B(t)| = 0$), lendo um valor inválido da memória.
\end{enumerate}

\subsection{Invariantes que a Solução deve Garantir}

A solução correta deve assegurar as seguintes propriedades:

\begin{center}
\renewcommand{\baselinestretch}{1.2}\normalsize
\begin{tabular}{|p{3.8cm}|p{9.5cm}|}
\hline
\textbf{Propriedade} & \textbf{Descrição} \\
\hline
Exclusão Mútua       & No máximo uma thread acessa o buffer por vez. \\
\hline
Ausência de Deadlock & O sistema sempre progride; nenhuma thread espera
                       indefinidamente por recurso que nunca será liberado. \\
\hline
Ausência de Starvation & Todo produtor eventualmente insere seu item; todo
                         consumidor eventualmente retira um item. \\
\hline
Consistência do Buffer & $0 \leq |B(t)| \leq N$ em qualquer instante. \\
\hline
\end{tabular}
\renewcommand{\baselinestretch}{1.5}\normalsize
\end{center}

% ===========================
%  3. DESCRICAO DO ALGORITMO
% ===========================
\section{Descrição do Algoritmo Desenvolvido}
\label{sec:algoritmo}

\subsection{Arquitetura do Sistema Refatorado}

A versão 2.0 adota uma arquitetura modular com responsabilidades bem separadas.
O projeto é composto pelos seguintes arquivos:

\begin{center}
\renewcommand{\baselinestretch}{1.2}\normalsize
\begin{tabular}{|p{3.6cm}|p{9.7cm}|}
\hline
\textbf{Arquivo} & \textbf{Responsabilidade} \\
\hline
\texttt{config.h}     & Constantes globais (\texttt{TAMANHO\_BUFFER}, \texttt{NUM\_PRODUTORES}, etc.) \\
\hline
\texttt{buffer.h/.c}  & Buffer circular + mutex + variáveis de condição \\
\hline
\texttt{logger.h/.c}  & Fila de log + semáforos POSIX + thread de log \\
\hline
\texttt{produtor.h/.c} & Lógica da thread produtora \\
\hline
\texttt{consumidor.h/.c} & Lógica da thread consumidora \\
\hline
\texttt{main.c}       & Inicialização, criação de threads e encerramento \\
\hline
\texttt{Makefile}     & Compilação de todos os módulos \\
\hline
\end{tabular}
\renewcommand{\baselinestretch}{1.5}\normalsize
\end{center}

\subsection{Visão Geral da Arquitetura}

A solução é composta por quatro camadas lógicas:

\begin{itemize}
    \item \textbf{Threads Produtoras} ($P_1, P_2, P_3$): geram itens e os
          inserem no buffer compartilhado; delegam o log via \texttt{log\_evento()}.
    \item \textbf{Buffer Circular Compartilhado} (5 posições): estrutura de dados
          central protegida por mutex e variáveis de condição.
    \item \textbf{Threads Consumidoras} ($C_1, C_2, C_3$): retiram e processam
          os itens do buffer; delegam o log via \texttt{log\_evento()}.
    \item \textbf{Thread de Logger Assíncrono}: consome eventos da fila de log
          e os exibe com timestamp; sincronizada por semáforos POSIX.
\end{itemize}

\subsection{Estruturas de Dados}

O buffer compartilhado é implementado como um \textbf{vetor circular} de
inteiros de tamanho fixo $N = 5$. Dois ponteiros — \texttt{cabeca} e
\texttt{cauda} — percorrem o vetor de forma circular usando aritmética modular,
evitando realocações dinâmicas. A variável \texttt{quantidade} mantém o número
de itens presentes no buffer em qualquer instante:

\begin{center}
\texttt{cabeca = (cabeca + 1) \% TAMANHO\_BUFFER} \\
\texttt{cauda  = (cauda  + 1) \% TAMANHO\_BUFFER}
\end{center}

A fila de log é também um vetor circular separado de \texttt{EventoLog}
com capacidade \texttt{TAMANHO\_FILA\_LOG = 20}.

\subsection{Mecanismos de Sincronização}

A tabela a seguir resume os três mecanismos utilizados na versão 2.0:

\begin{center}
\renewcommand{\baselinestretch}{1.2}\normalsize
\begin{tabular}{|p{4.0cm}|p{3.5cm}|p{5.8cm}|}
\hline
\textbf{Primitiva} & \textbf{Onde é usada} & \textbf{Finalidade} \\
\hline
\texttt{pthread\_mutex\_t} & Buffer principal & Exclusão mútua na região crítica de inserção/remoção \\
\hline
\texttt{pthread\_cond\_t} & Buffer principal & Bloquear produtores (buffer cheio) e consumidores (buffer vazio) \\
\hline
\texttt{sem\_t} (POSIX) & Fila de log & Controlar espaços livres e itens pendentes na fila de eventos \\
\hline
\end{tabular}
\renewcommand{\baselinestretch}{1.5}\normalsize
\end{center}

\subsection{Protocolo da Thread Produtora}

Cada thread produtora segue o protocolo abaixo:

\begin{enumerate}
    \item Simular tempo de produção (\texttt{usleep}).
    \item Invocar \texttt{buffer\_inserir(id)}, que internamente:
    \begin{enumerate}
        \item Adquire o mutex (\texttt{pthread\_mutex\_lock}).
        \item \textbf{Enquanto} o buffer estiver cheio, envia \texttt{LOG\_BUF\_CHEIO}
              ao logger e aguarda na variável de condição \texttt{nao\_cheio}.
        \item Insere o item no slot \texttt{cauda} e avança o ponteiro.
        \item Incrementa \texttt{quantidade}.
        \item Sinaliza a variável \texttt{nao\_vazio}.
        \item Libera o mutex.
        \item Envia evento \texttt{LOG\_PRODUZIU} ao logger assíncrono.
    \end{enumerate}
\end{enumerate}

O uso do laço \texttt{while} ao invés de \texttt{if} na etapa 2b é essencial
para lidar com \textit{spurious wakeups} — situação prevista pelo padrão POSIX
em que a thread pode ser desbloqueada sem que a condição esperada tenha se
tornado verdadeira \cite{kerrisk}.

\subsection{Logger Assíncrono com Semáforos POSIX}

O módulo de log (\texttt{logger.c}) implementa uma fila produtora-consumidora
própria, onde \textit{produtores são as threads do sistema principal} e o
\textit{único consumidor é a thread de log}. A sincronização usa
\textbf{exclusivamente semáforos POSIX} (\texttt{sem\_t}):

\begin{itemize}
    \item \texttt{sem\_t sem\_espacos\_log}: semáforo contador inicializado com
          \texttt{TAMANHO\_FILA\_LOG} (20). Representa espaços livres na fila.
    \item \texttt{sem\_t sem\_itens\_log}: semáforo contador inicializado com 0.
          Representa itens pendentes para exibição.
    \item \texttt{pthread\_mutex\_t mutex\_log}: protege o índice circular da
          fila de log durante leitura e escrita.
\end{itemize}

O fluxo de envio de um evento via \texttt{log\_evento()} é:

\begin{enumerate}
    \item Chama \texttt{sem\_trywait(\&sem\_espacos\_log)} --- não bloqueia.
    \item Se retornar \texttt{EAGAIN} (fila cheia), descarta o evento silenciosamente.
    \item Caso contrário, preenche o evento com timestamp e escreve na fila com proteção do \texttt{mutex\_log}.
    \item Chama \texttt{sem\_post(\&sem\_itens\_log)} para sinalizar a thread de log.
\end{enumerate}

A thread de log executa em loop:

\begin{enumerate}
    \item \texttt{sem\_wait(\&sem\_itens\_log)} --- bloqueia enquanto fila vazia.
    \item Adquire \texttt{mutex\_log}, lê o evento, libera o mutex.
    \item \texttt{sem\_post(\&sem\_espacos\_log)} --- libera uma posição.
    \item Exibe o evento formatado com timestamp em milissegundos.
    \item Para ao receber evento sentinela \texttt{LOG\_SHUTDOWN}.
\end{enumerate}

O uso de \texttt{sem\_trywait} (não-bloqueante) no envio garante que
produtores e consumidores \textbf{nunca bloqueiam} aguardando o logger,
preservando o desempenho do sistema principal.

\subsection{Protocolo da Thread Consumidora}

A thread consumidora segue o protocolo simétrico ao produtor:

\begin{enumerate}
    \item Simular tempo de consumo (\texttt{usleep}).
    \item Invocar \texttt{buffer\_retirar(id)}, que internamente:
    \begin{enumerate}
        \item Adquire o mutex.
        \item \textbf{Enquanto} o buffer estiver vazio, envia \texttt{LOG\_BUF\_VAZIO}
              e aguarda na variável \texttt{nao\_vazio}.
        \item Lê o item no slot \texttt{cabeca} e avança o ponteiro.
        \item Decrementa \texttt{quantidade}.
        \item Sinaliza \texttt{nao\_cheio}.
        \item Libera o mutex.
        \item Envia evento \texttt{LOG\_CONSUMIU} ao logger assíncrono.
    \end{enumerate}
\end{enumerate}

\subsection{Correção da Solução}

A tabela a seguir mapeia cada propriedade de correção exigida ao mecanismo
que a garante na implementação:

\begin{center}
\renewcommand{\baselinestretch}{1.2}\normalsize
\begin{tabular}{|p{3.8cm}|p{9.5cm}|}
\hline
\textbf{Propriedade} & \textbf{Mecanismo} \\
\hline
Exclusão Mútua
    & \texttt{pthread\_mutex\_lock/unlock} envolve toda a região crítica. \\
\hline
Ausência de Overflow
    & Laço \texttt{while (quantidade == N)} bloqueia o produtor antes de inserir. \\
\hline
Ausência de Underflow
    & Laço \texttt{while (quantidade == 0)} bloqueia o consumidor antes de retirar. \\
\hline
Ausência de Deadlock
    & O mutex é liberado por \texttt{pthread\_cond\_wait} e os sinais
      são emitidos após cada inserção/remoção. \\
\hline
Ausência de Starvation
    & O escalonador POSIX garante progresso; \texttt{pthread\_cond\_signal}
      acorda pelo menos uma thread em espera. \\
\hline
Log não-bloqueante
    & \texttt{sem\_trywait} descarta eventos se a fila de log estiver cheia,
      sem atrasar produtores nem consumidores. \\
\hline
\end{tabular}
\renewcommand{\baselinestretch}{1.5}\normalsize
\end{center}

\subsection{Resultado da Execução}

A execução do programa com 3 produtores, 3 consumidores e buffer de 5 posições
produziu a saída abaixo, confirmando a completa sincronização e ausência de
condições de corrida. O timestamp em milissegundos evidencia o caráter
assíncrono do logger:

\begin{quote}
{\ttfamily\small\renewcommand{\baselinestretch}{1.1}\normalsize
=== Produtor-Consumidor v2.0 com Logger Assincrono ===\\
Buffer: 5 posicoes | 3 produtores | 3 consumidores\\[2pt]
{[}  158.6ms{]} {[}Produtor 3{]} Inseriu: 1 | Buffer: 1/5\\
{[}  212.4ms{]} {[}Produtor 2{]} Inseriu: 2 | Buffer: 2/5\\
{[}  241.5ms{]} {[}Produtor 1{]} Inseriu: 3 | Buffer: 3/5\\
{[}  281.4ms{]} {[}Consumidor 3{]} Retirou: 1 | Buffer: 2/5\\
{[}  366.4ms{]} {[}Consumidor 1{]} Retirou: 2 | Buffer: 1/5\\
{[}  376.8ms{]} {[}Produtor 1{]} Inseriu: 4 | Buffer: 2/5\\
{[}  379.8ms{]} {[}Produtor 3{]} Inseriu: 5 | Buffer: 3/5\\
{[}  540.5ms{]} {[}Consumidor 2{]} Retirou: 3 | Buffer: 2/5\\
{[}  552.5ms{]} {[}Produtor 2{]} Inseriu: 6 | Buffer: 3/5\\
{[}  596.9ms{]} {[}Produtor 3{]} Inseriu: 7 | Buffer: 4/5\\
{[}  620.0ms{]} {[}Produtor 1{]} Inseriu: 8 | Buffer: 5/5\\
{[}  624.6ms{]} {[}Consumidor 3{]} Retirou: 4 | Buffer: 4/5\\
{[}  748.7ms{]} {[}Produtor 2{]} Inseriu: 9 | Buffer: 5/5\\
{[}  759.1ms{]} {[}Produtor 3{]} Buffer cheio. Aguardando...\\
{[}  840.7ms{]} {[}Consumidor 1{]} Retirou: 5 | Buffer: 4/5\\
{[}  840.8ms{]} {[}Produtor 3{]} Inseriu: 10 | Buffer: 5/5\\
{[}  844.7ms{]} {[}Consumidor 2{]} Retirou: 6 | Buffer: 4/5\\
{[} 1012.3ms{]} {[}Produtor 1{]} Inseriu: 11 | Buffer: 5/5\\
{[} 1046.8ms{]} {[}Produtor 2{]} Buffer cheio. Aguardando...\\
{[} 1072.7ms{]} {[}Consumidor 3{]} Retirou: 7 | Buffer: 4/5\\
{[} 1072.8ms{]} {[}Produtor 2{]} Inseriu: 12 | Buffer: 5/5\\
{[} 1099.9ms{]} {[}Consumidor 1{]} Retirou: 8 | Buffer: 4/5\\
{[} 1312.1ms{]} {[}Consumidor 2{]} Retirou: 9 | Buffer: 3/5\\
{[} 1509.1ms{]} {[}Consumidor 1{]} Retirou: 10 | Buffer: 2/5\\
{[} 1625.0ms{]} {[}Consumidor 3{]} Retirou: 11 | Buffer: 1/5\\
{[} 1745.4ms{]} {[}Consumidor 2{]} Retirou: 12 | Buffer: 0/5\\[2pt]
=== Execucao concluida com sucesso! ===
}
\end{quote}

A análise da saída demonstra precisamente as seguintes propriedades:
\begin{enumerate}
    \item \textbf{Bloqueio em buffer cheio:} em 759,1\,ms e 1046,8\,ms, os produtores
          encontraram o buffer cheio (5/5) e suspenderam sua execução corretamente
          via variável de condição \texttt{nao\_cheio}.
    \item \textbf{Retomada pontual:} logo que um consumidor liberou vaga
          (em 840,7\,ms e 1072,7\,ms), o produtor suspenso acordou prontamente
          e concluiu a inserção.
    \item \textbf{Consumo integral e esvaziamento:} exatamente 12 itens foram
          produzidos (3 produtores $\times$ 4 itens) e todos os 12 itens foram
          consumidos (3 consumidores $\times$ 4 itens), resultando na ocupação final
          de 0/5 no instante 1745,4\,ms.
    \item \textbf{Sequência correta de encerramento:} a thread principal
          aguardou o término de todos os produtores e em seguida de todos os
          consumidores antes de invocar \texttt{logger\_shutdown()}, garantindo
          que nenhum registro de log fosse perdido ou truncado.
\end{enumerate}

% ===========================
%  4. CONCLUSAO
% ===========================
\section{Conclusão}
\label{sec:conclusao}

Este trabalho, em sua Versão 2.0, expandiu a solução original do Problema
Produtor-Consumidor em duas dimensões fundamentais: a \textbf{refatoração
modular} do código e a \textbf{adição de semáforos POSIX} como nova primitiva
de sincronização.

A refatoração separou claramente as responsabilidades em módulos independentes
(\texttt{buffer}, \texttt{logger}, \texttt{produtor}, \texttt{consumidor},
\texttt{main}), facilitando a manutenção, o teste e a compreensão do sistema.
O Makefile automatiza a compilação de todos os módulos.

O logger assíncrono demonstrou a aplicação prática de \textbf{semáforos POSIX}
(\texttt{sem\_t}) como alternativa às variáveis de condição: enquanto o buffer
principal emprega o par mutex e condição, a fila de log combina os semáforos
\texttt{sem\_espacos\_log} e \texttt{sem\_itens\_log} com o lock \texttt{mutex\_log}.
O uso de \texttt{sem\_trywait} garantiu que as threads do sistema principal nunca
bloqueiem aguardando o logger.

Como lições aprendidas, destacam-se:

\begin{itemize}
    \item O uso do laço \texttt{while} ao invés de \texttt{if} ao aguardar
          uma variável de condição é indispensável para robustez frente a
          \textit{spurious wakeups}.
    \item O padrão \textit{lock~→~verificar~→~operar~→~sinalizar~→~unlock}
          deve ser aplicado de forma consistente para evitar janelas de
          vulnerabilidade.
    \item Semáforos POSIX são especialmente adequados para padrões
          produtor-consumidor simples, pois encapsulam diretamente o contador
          de itens e espaços sem a necessidade de um mutex separado para a
          condição.
    \item A combinação de múltiplas primitivas (mutex, variáveis de condição
          e semáforos) em um mesmo sistema requer atenção à ordem de
          encerramento para evitar deadlocks.
\end{itemize}

Como trabalhos futuros, sugere-se explorar o problema dos \textit{Leitores-Escritores}
e o \textit{Jantar dos Filósofos} com as mesmas primitivas, bem como a comparação
de desempenho entre mutex+condição e semáforos POSIX em cenários de alta contention.

% ===========================
%  REFERENCIAS
% ===========================
\newpage
\renewcommand{\refname}{Referências}
\begin{thebibliography}{9}

\bibitem{tanenbaum}
TANENBAUM, A. S.; BOS, H.
\textbf{Modern Operating Systems}.
4.~ed. Upper Saddle River: Pearson, 2015.

\bibitem{silberschatz}
SILBERSCHATZ, A.; GALVIN, P. B.; GAGNE, G.
\textbf{Operating System Concepts}.
10.~ed. Hoboken: Wiley, 2018.

\bibitem{kerrisk}
KERRISK, M.
\textbf{The Linux Programming Interface}.
San Francisco: No Starch Press, 2010.

\bibitem{dijkstra}
DIJKSTRA, E. W.
Cooperating Sequential Processes.
In: \textit{Programming Languages}, F. Genuys (ed.),
Academic Press, New York, 1968, pp.~43--112.

\bibitem{buttazzo}
BUTTAZZO, G. C.
\textbf{Hard Real-Time Computing Systems}.
3.~ed. New York: Springer, 2011.

\bibitem{pthread_man}
THE OPEN GROUP.
\textbf{POSIX.1-2017 -- pthread.h}.
Disponível em: \url{https://pubs.opengroup.org/onlinepubs/9699919799/}.
Acesso em: nov.~2026.

\end{thebibliography}

% ===========================
%  APENDICE: Codigo-Fonte
% ===========================
\newpage
\appendix
\section{Código-Fonte dos Módulos Principais}

\subsection*{config.h}

{\ttfamily\small\renewcommand{\baselinestretch}{1.1}\normalsize
\noindent
\#ifndef CONFIG\_H\\
\#define CONFIG\_H\\[4pt]
\#define TAMANHO\_BUFFER\ \ \ 5\\
\#define NUM\_PRODUTORES\ \ \ 3\\
\#define NUM\_CONSUMIDORES\ 3\\
\#define ITENS\_POR\_THREAD\ 4\\
\#define TAMANHO\_FILA\_LOG\ 20\\[4pt]
\#endif
}

\subsection*{buffer.h}

{\ttfamily\small\renewcommand{\baselinestretch}{1.1}\normalsize
\noindent
\#ifndef BUFFER\_H\\
\#define BUFFER\_H\\[4pt]
\#include $<$pthread.h$>$\\
\#include "config.h"\\[4pt]
typedef struct \{\\
\hspace*{1em}int dados[TAMANHO\_BUFFER];\\
\hspace*{1em}int cabeca, cauda, quantidade;\\
\hspace*{1em}int proximo\_valor;\\
\} Buffer;\\[4pt]
extern Buffer buffer\_compartilhado;\\
extern pthread\_mutex\_t mutex\_buffer;\\
extern pthread\_cond\_t\ \ nao\_cheio;\\
extern pthread\_cond\_t\ \ nao\_vazio;\\[4pt]
void buffer\_init(void);\\
void buffer\_destroy(void);\\
int\ \ buffer\_inserir(int id\_thread);\\
int\ \ buffer\_retirar(int id\_thread);\\[4pt]
\#endif
}

\subsection*{logger.h}

{\ttfamily\small\renewcommand{\baselinestretch}{1.1}\normalsize
\noindent
\#ifndef LOGGER\_H\\
\#define LOGGER\_H\\[4pt]
typedef enum \{\\
\hspace*{1em}LOG\_PRODUZIU, LOG\_CONSUMIU,\\
\hspace*{1em}LOG\_BUF\_CHEIO, LOG\_BUF\_VAZIO,\\
\hspace*{1em}LOG\_SHUTDOWN\\
\} TipoEvento;\\[4pt]
typedef struct \{\\
\hspace*{1em}TipoEvento tipo;\\
\hspace*{1em}int\ \ thread\_id;\\
\hspace*{1em}int\ \ valor;\\
\hspace*{1em}int\ \ quantidade\_no\_buffer;\\
\hspace*{1em}double timestamp\_ms;\\
\} EventoLog;\\[4pt]
void logger\_init(void);\\
void logger\_iniciar\_thread(void);\\
void log\_evento(TipoEvento, int, int, int);\\
void logger\_shutdown(void);\\
void logger\_destroy(void);\\[4pt]
\#endif
}

\subsection*{logger.c (trechos principais)}

{\ttfamily\small\renewcommand{\baselinestretch}{1.1}\normalsize
\noindent
/* Semaforos POSIX */\\
static sem\_t sem\_espacos\_log;\ /* espacos livres */\\
static sem\_t sem\_itens\_log;\ \ \ /* itens pendentes */\\
static pthread\_mutex\_t mutex\_log = PTHREAD\_MUTEX\_INITIALIZER;\\[4pt]
void logger\_init(void) \{\\
\hspace*{1em}clock\_gettime(CLOCK\_MONOTONIC, \&tempo\_inicio);\\
\hspace*{1em}sem\_init(\&sem\_espacos\_log, 0, TAMANHO\_FILA\_LOG);\\
\hspace*{1em}sem\_init(\&sem\_itens\_log,\ \ 0, 0);\\
\}\\[4pt]
void log\_evento(TipoEvento tipo, int id, int val, int qtd) \{\\
\hspace*{1em}if (sem\_trywait(\&sem\_espacos\_log) != 0) return; /* descarta */\\
\hspace*{1em}EventoLog ev = \{tipo, id, val, qtd, ms\_desde\_inicio()\};\\
\hspace*{1em}pthread\_mutex\_lock(\&mutex\_log);\\
\hspace*{1em}fila\_log[cauda\_log] = ev;\\
\hspace*{1em}cauda\_log = (cauda\_log + 1) \% TAMANHO\_FILA\_LOG;\\
\hspace*{1em}pthread\_mutex\_unlock(\&mutex\_log);\\
\hspace*{1em}sem\_post(\&sem\_itens\_log);\\
\}\\[4pt]
static void *thread\_log\_func(void *arg) \{\\
\hspace*{1em}while (1) \{\\
\hspace*{2em}sem\_wait(\&sem\_itens\_log);\\
\hspace*{2em}pthread\_mutex\_lock(\&mutex\_log);\\
\hspace*{2em}EventoLog ev = fila\_log[cabeca\_log];\\
\hspace*{2em}cabeca\_log = (cabeca\_log + 1) \% TAMANHO\_FILA\_LOG;\\
\hspace*{2em}pthread\_mutex\_unlock(\&mutex\_log);\\
\hspace*{2em}sem\_post(\&sem\_espacos\_log);\\
\hspace*{2em}if (ev.tipo == LOG\_SHUTDOWN) break;\\
\hspace*{2em}/* exibe o evento com timestamp */\\
\hspace*{1em}\}\\
\hspace*{1em}return NULL;\\
\}
}

\subsection*{main.c (estrutura de terminação)}

{\ttfamily\small\renewcommand{\baselinestretch}{1.1}\normalsize
\noindent
int main(void) \{\\
\hspace*{1em}buffer\_init();\ logger\_init();\ logger\_iniciar\_thread();\\[4pt]
\hspace*{1em}/* Cria consumidoras primeiro, depois produtoras */\\
\hspace*{1em}for (i = 0; i $<$ NUM\_CONSUMIDORES; i++)\\
\hspace*{2em}pthread\_create(\&tc[i], NULL, consumidor, \&idc[i]);\\
\hspace*{1em}for (i = 0; i $<$ NUM\_PRODUTORES; i++)\\
\hspace*{2em}pthread\_create(\&tp[i], NULL, produtor, \&idp[i]);\\[4pt]
\hspace*{1em}/* Join produtores primeiro */\\
\hspace*{1em}for (i = 0; i $<$ NUM\_PRODUTORES; i++)\\
\hspace*{2em}pthread\_join(tp[i], NULL);\\[4pt]
\hspace*{1em}/* Join consumidores (garante consumo completo do buffer) */\\
\hspace*{1em}for (i = 0; i $<$ NUM\_CONSUMIDORES; i++)\\
\hspace*{2em}pthread\_join(tc[i], NULL);\\[4pt]
\hspace*{1em}/* Encerra o logger (flush da fila apos todo processamento) */\\
\hspace*{1em}logger\_shutdown();\\[4pt]
\hspace*{1em}buffer\_destroy();\ logger\_destroy();\\
\hspace*{1em}return 0;\\
\}
}

\end{document}
